\section{Planificación de actividades}

Las actividades planificadas y realizadas durante el periodo de residencias enero-junio del 2026 son las que se describen a continuación. En la figura \ref{fig:Cronograma} se observa el tiempo designado para cada fase y un avance real hasta la entrega de este reporte.

Fase 1: Familiarización con el entorno, el SIP y el stack tecnológico
\begin{itemize}
    \item Inducción al área y revisión de procesos internos.
    \item Exploración de la arquitectura general del SIP y sus módulos principales.
\end{itemize}

Fase 2: Prácticas iniciales de React y construcción de componentes básicos
\begin{itemize}
    \item Repaso de conceptos esenciales: hooks, estados, props.
    \item Construcción de componentes simples.
\end{itemize}

Fase 3: Implementación de componentes reutilizables orientados al SIP
\begin{itemize}
    \item Identificación de componentes candidatos: tablas, formularios, modales, filtros.
    \item Integración de estilos estandarizados (temas, tipografías, espaciados).
\end{itemize}

Fase 4: Desarrollo de vistas funcionales dentro del SIP
\begin{itemize}
    \item Maquetado inicial de pantallas usando componentes creados.
    \item Implementación de lógica de negocio local (validaciones, comportamientos).
    \item Integración de componentes reutilizables dentro del flujo real de las vistas.
\end{itemize}

Fase 5: Integración con servicios y manejo de datos reales o simulados
\begin{itemize}
    \item Consumo de APIs: fetch, manejo de errores.
    \item Pruebas de conexión en entornos de desarrollo.
    \item Verificación con datos reales y retroalimentación con el área funcional.
\end{itemize}

Fase 6: Ajustes, pruebas, mejoras visuales y refinamiento general
\begin{itemize}
    \item Corrección de errores detectados durante QA interno.
    \item Mejoras visuales.
\end{itemize}

Fase 7: Documentación y presentación de avances del proyecto
\begin{itemize}
    \item Elaboración de documentación técnica final (componentes, vistas, APIs usadas).x
    \item Demostración del proyecto final.
\end{itemize}

Estas actividades aseguran la entrega de una solución tecnológica robusta, moderna y alineada con las necesidades operativas de la organización.
\begin{figure}[h]
    \includegraphics[width=1\linewidth]{cronograma.png}
    \caption{Cronograma de trabajo.}
    \label{fig:Cronograma}
\end{figure}
